\documentclass[a4j,8pt,fleqn]{jarticle}
\usepackage[dvipdfm]{graphicx,color}
\usepackage{graphicx,layout,url,comment,here,listings,here,ascmac,amsmath,amssymb}
\usepackage{./sty/algorithm,./sty/algorithmic,./sty/framed,./sty/jlisting,./sty/doxygenorig,./sty/eclepsf,./sty/listliketab,./sty/mediabb,./sty/misc,./sty/ncodeline,./sty/page4prop,./sty/subfigure}
\usepackage{./sty/engset,./sty/style}

% 余白の設定
% \usepackage[top=30truemm,bottom=30truemm,left=25truemm,right=25truemm]{geometry}

% ソースコードのシンタックスハイライト
\lstset{language=c,
  breaklines=true,
  basicstyle=\ttfamily\small,
  commentstyle={\small\itshape \color[cmyk]{1,0.4,1,0}},
  classoffset=1,
  keywordstyle={\small\bfseries \color[cmyk]{0,1,0,0}},
  stringstyle={\small\ttfamily \color[rgb]{0,0,1}},
  frame=tblr,
  framesep=5pt,
  showstringspaces=false,
  numbers=left,
  stepnumber=1,
  numberstyle=\tiny,
  tabsize=2,
  xrightmargin=2zw,
  xleftmargin=2zw,
  numbersep=2zw,
}

\makeatletter

% ソースコードの改行
\def\lst@lettertrue{\let\lst@ifletter\iffalse}

% 参考文献リストのサイズ
\renewenvironment{thebibliography}[1]
{\section*{\refname\@mkboth{\refname}{\refname}}%
  \list{\@biblabel{\@arabic\c@enumiv}}%
       {\settowidth\labelwidth{\@biblabel{#1}}%
        \leftmargin\labelwidth
        \advance\leftmargin\labelsep
 \setlength\itemsep{-0.5zh}%
 \setlength\baselineskip{0.5pt}%
        \@openbib@code
        \usecounter{enumiv}%
        \let\p@enumiv\@empty
        \renewcommand\theenumiv{\@arabic\c@enumiv}}%
  \sloppy
  \clubpenalty4000
  \@clubpenalty\clubpenalty
  \widowpenalty4000%
  \sfcode`\.\@m}
 {\def\@noitemerr
   {\@latex@warning{Empty `thebibliography' environment}}%
  \endlist}

\makeatother

% アルゴリズム表示
\renewcommand{\algorithmicrequire}{\textbf{Input:}}
\renewcommand{\algorithmicensure}{\textbf{Output:}}

\pagestyle{empty}

\begin{document}

% 和文タイトル
\jtitle{タイトル}

% 英文タイトル
\ifENG
\etitle{Title}
\fi

% 和文著者
\jauthor{名前\dag}

% 英文著者
\ifENG
\eauthor{Name\dag}
\fi

% 和文連絡先
\jcontact{\dag 大学 学部\\000-0000 住所\\\texttt{mail@mail.com}\\[1ex]}

% 英文連絡先
\ifENG
\econtact{\dag Faculty, University\\Address\\\texttt{mail@mail.com}\\[1ex]}
\fi

% 和文アブスト
\begin{jabstract}
これは忘れた頃にソースコードを見てLaTeXの書き方を思い出すためのメモ書きです．
\end{jabstract}

% 英文アブスト
\ifENG
\begin{eabstract}
Abstract
\end{eabstract}
\fi

\maketitle
\thispagestyle{empty}

\section{基本書式}
こんにちは．引用\cite{sample}してみた．\par
改行もしてみた\footnote{なお注釈はこのように書ける}．

\subsection{箇条書き}
これは普通の箇条書き．
\begin{itemize}
\item あああ
\end{itemize}
\par
これは説明付き箇条書き．
\begin{description}
 \item[仮説A]あああ
 \item[仮説B]いいい
\end{description}
\par
これは番号付き箇条書き．
\begin{enumerate}
  \item あああ
  \item いいい
\end{enumerate}

\subsection{図}
図\ref{fig:one}を挿入してみた．
\begin{figure}[H]
  \begin{center}
    \includegraphics[width=20mm]{./img/sample.png}
    \small
    \caption{図のサンプル}
    \label{fig:one}
  \end{center}
\end{figure}
\par
図を入れるには画像を変換すること．
\begin{framed}
  \begin{small}
    \begin{verbatim}
ebb sample.png
    \end{verbatim}
  \end{small}
\end{framed}
\par
あと，次のように設定するとその場に図表やソースコードを描画できる．
\begin{framed}
  \begin{small}
    \begin{verbatim}
\usepackage{here}
...
\begin{figure}[H]
    \end{verbatim}
  \end{small}
\end{framed}

\subsection{表}
表はこんな感じ．
\begin{table}[htb]
  \begin{tabular}{|l|c|c|}\hline
    あああ & いいい & ううう \\\hline
    えええ & おおお & かかか\\\hline
  \end{tabular}
\end{table}

\subsection{数式}
試しにFitzHugh-Nagumo modelの式を書いておく．ここで\\
\begin{equation}
  \dfrac {dv} {dt} = \dfrac {1}{ε} (v - \dfrac {v^3}{3} -w) + I_{ext}\\
\end{equation}
\begin{equation}
  \dfrac {dw} {dt} = ε(v - β - γw) 
\end{equation}
ただし$v$は細胞の膜電位，$w$は細胞の不活性化の程度を表す抽象的な変数，$I_{ext}$は細胞外から届く電流を表す．

\newpage % 改ページ

\subsection{アルゴリズム}
アルゴリズム\ref{alg:one}は次のように記述する．
\begin{algorithm}[H]
  \begin{small}
    \caption{アルゴリズム}
    \label{alg:one}
    \begin{algorithmic}
\REQUIRE 入力
\ENSURE 出力
\STATE 文
\WHILE {true}
\STATE 文
\ENDWHILE
\IF {条件}
\STATE {文}
\ENDIF
    \end{algorithmic}
  \end{small}
\end{algorithm}

\subsection{ソースコード}
ソースコードは次のように記述する．
\begin{lstlisting}[caption=hello.c][H]
// コメント
#include <stdio.h>

int main(int argc, char *argv[])
{
  return 0;
}
\end{lstlisting}

\begin{thebibliography}{20}
%\bibitem{a}
\bibitem{sample}sample\\\texttt{https://localhost/}
\end{thebibliography}
\end{document}

% EOF